\section{Introducción}
    \subsection{¿Que son los numeros amigos?}
    Dos números son amigos cuando:
        - la suma de los divisores del primero (excepto él mismo) es igual al segundo
        - la suma de los divisores del segundo (excepto él mismo) es igual al primero

        \subsection{Ejemplo}
            Los números 220 y 284 son amigos: 
            
            Los divisores de 220 son 1, 2, 4, 5, 10, 11, 20, 22, 44, 55 y 110, y su suma es 284 

            Los divisores de 284 son 1, 2, 4, 71 y 142, que suman 220

        \subsection{¿Como resolverlo?}
            La forma mas directa de resolverlo es haciendo un algoritmo que recorra todos los números hasta el valor máximo dado, y para cada uno de ellos calcule la suma de sus divisores (excepto él mismo) y verifique si ese resultado
            tiene como suma de sus divisores al numero por el que está iterando.
            La consecuencia que tiene este algoritmo es que su complejidad temporal es muy alta ya que recorre y recalcula divisores y sumas de divisores muchas veces.
            
        \subsection{Supuestos}
            Se supone como valor maximo a hallar siempre un numero entero positivo, mayor o igual a 0 y como salida del algoritmo no un listado de tuplas de pares de amigos
            sino la impresion por consola de esta, por lo que no hace falta estructuras de datos internas del algoritmo que requieran de almacenar estos encontrados para 
            poder imprimirlos
            
        \subsection{Implementación}
            \begin{lstlisting}[language=Python]
import time
def amigos(MAX):
    
    t1=time.time()
    for i in range(MAX):
        s=0
        for j in range (1,i-1+1):
            if i%j==0:
                s+=j
        s2=0
        for k in range (1,s-1):
            if s%k==0:
                s2+=k
        if i==s2:
            print(i,s)
   
    t2=time.time()    
    print(t2-t1)
    
amigos(100000)
            \end{lstlisting}
        \subsection{Consecuencia}
            El algoritmo tarda 329 segundos en encontrar los pares de numeros amigos hasta el 100.000 (en el ordenador del escritor del informe)\footnote{Procesador: Intel Core I5, single thread, 4.3GHz max} y 
            en el siguiente grafico se visualiza el crecimiento de los tiempos a medida que aumenta el valor maximo hasta el cual se busca.